%% kapa: exam=1 round=35 year=2024 period=1 form=A subject=civil_law subject_ko=민법 seq=10 number=10
%% source: Q-Net 감정평가사 기출문제 (cst00309) artlSeq=5213229 file=raw/1008/5213229/1교시 문제지 원본.pdf
%% conversion: kordoc 4.24.51 -> tools/gen1.py (auto); review=proofread-sonnet; flags=none
\begin{question}{10}
의사표시에 관한 설명으로 옳지 않은 것은? (다툼이 있으면 판례에 따름)
\begin{choices}
  \item 의사표시자가 통지를 발송한 후 사망하더라도 그 의사표시의 효력에 영향을 미치지 않는다.
  \item 통정허위표시의 경우, 통정의 동기나 목적은 허위표시의 성립에 영향이 없다.
  \item 통정허위표시로 무효인 경우, 당사자는 가장행위의 채무불이행이 있더라도 이를 이유로 하는 손해배상을 청구할 수 없다.
  \item 착오로 인하여 표의자가 경제적 불이익을 입지 않는 경우에는 특별한 사정이 없는 한 중요부분의 착오라고 할 수 없다.
  \item 상대방이 표의자의 착오를 알고 이용하였더라도 착오가 표의자의 중대한 과실로 인한 경우에는 표의자는 착오를 이유로 그 의사표시를 취소할 수 없다.
\end{choices}
\end{question}
\answer{5}
